\scanpage{23}
% L1–3
On a ainsi prouvé le

\uline{Théorème} $\unsure{H\!H|H^1}(A,B)$ s'identifie \deleted{à $A$} au groupe abélien
\par\noindent\begin{tabular}{|p{.92\linewidth}}
$A/A^p$ (où $A^p$ est l'ens des $x^p$, avec $x\in A$).
% L4–7
Si \deleted{\illegible{1}} $a\in A^{\deleted{\unsure{p}}}$, il lui correspond la classe de la $A$-algèbre $C_a=A[U]/(U^p-a)$, \unsure{muni|munie} de la $A$-dérivation $D_a$ définie par $D_aU=1$, et $C_a\simeq C_b$ si et \unsure{seult|seulement} si $a\equiv b$ $(A^p)$.

% L8–12
Soit $C$ un $B$-espace homogène principal, alors il existe dans $C$ un $u$ \deleted{\unsure{$\in A$|et}} tel que $Du=1$, et \deleted{\unsure{avec $a\in{}$ alors}} on a $u^p=a\in A$, $C\simeq C_a$ en \deleted{tant} que \deleted{fibré} $B$-fibré\,; $u$ est bien déterminé mod addition d'un élément de $A$.
\end{tabular}

% L13–15
On n'a pas prouvé que la loi additive naturelle de $H\!H(A;B)$ correspond à celle de $A/A^p$, mais ça ne peut être que facile.

% L16–19
\uline{Corollaire} Soit $C$ une $A$-algèbre commutative. Pour
\par\noindent\begin{tabular}{|p{.92\linewidth}}
\deleted{sur laquelle} \inserted{qu'il} existe \deleted{une} sur $C$ une structure de fibré \inserted{hom.} principal sous $B$, il faut et il suffit que $C$ ait une $p$-base sur $A$ réduite à un élément, \unsure{eh|et} i.e. un $u\in C$
% L20–22
tel que les $u^i$ $(0\leq i\leq p-1)$ forment une base de $C$ sur $A$. La donnée d'une telle structure de $B$-fibré \inserted{hom.} principal sur $C$
\end{tabular}
% archive sheet number: 22
